
\paragraph{An order-saturation theorem for the division-free budget.}
The Duhamel condition avoids quotients that become
singular at multiple zeros. It does not, however,
admit an order-independent margin even at the
central Fourier frequency. The following result
is obtained directly from the genuine theta
heat kernel and does not assume RH.

\begin{theorem}[High-order saturation of the central heat budget]
\label{thm:duhamelsaturation}
Write $t_0=1/2$, let $0<s\le t_0$ be fixed, and put
\[
B_n(s)=(P_sU_n(0,\cdot))(0),\qquad
D_n(s)=k_n\int_0^s
(P_{s-r}U_{n+1}(r,\cdot))(0)\,dr,
\]
where $k_n=(n+1)(2n+1)$ and $P_s$
is the Gaussian semigroup for
$\frac12\partial_x^2$.
For the classical Jacobi-theta heat flow,
$B_n(s)>0$ for all $n\ge0$ and
\begin{equation}\label{eq:duhamelcentralratio}
0<\frac{U_n(s,0)}{B_n(s)}<1,\qquad
\frac{U_{n+1}(s,0)}{B_{n+1}(s)}
<
\frac{U_n(s,0)}{B_n(s)}.
\end{equation}
Moreover
\begin{equation}\label{eq:duhamelsaturation}
\boxed{\lim_{n\to\infty}
\frac{D_n(s)}{B_n(s)}=1.}
\end{equation}
In particular, for each fixed $s>0$ and
$\eta>0$ the strengthened requirement
$D_n(s)\le(1-\eta)B_n(s)$ fails for
all sufficiently large Laguerre orders.
\end{theorem}

\begin{proof}
Let
$\nu(u)=\frac12e^{t_0u^2}\Phi(u)>0$,
so that $H_{t_0}(x)=\int_{\mathbb R}
\nu(u)e^{ixu}\,du$.
Write $w=u+v$, $b=u-v$ and set
\[
W_s(b)=\frac12\int_{\mathbb R}
e^{-sw^2/2}
\nu\left(\frac{w+b}{2}\right)
\nu\left(\frac{w-b}{2}\right)\,dw.
\]
The density $W_s$ is continuous,
strictly positive on every real $b$,
and has moments of every order,
since the theta kernel decays faster
than any Gaussian on the real axis.
The Fourier correlation formula and
the Gaussian multiplier in $P_s$ give
\begin{align}
B_n(s)&=\frac1{(2n)!}
\int_{\mathbb R}b^{2n}W_s(b)\,db,
\label{eq:duhamelbn}\\
U_n(s,0)&=\frac1{(2n)!}
\int_{\mathbb R}b^{2n}e^{-sb^2/2}
W_s(b)\,db.
\label{eq:duhamelun}
\end{align}
Indeed, $u^2+v^2=(w^2+b^2)/2$.
Duhamel's identity gives
$D_n(s)=B_n(s)-U_n(s,0)$;
the integrand of the difference is
strictly positive for $b\ne0$.

To prove strict monotonicity, define
the probability measure
$d\mu_n(b)=b^{2n}W_s(b)\,db/
\int b^{2n}W_s(b)\,db$.
Then $r_n:=U_n(s,0)/B_n(s)
=\int e^{-sb^2/2}\,d\mu_n$.
The measure $\mu_{n+1}$ is the
$b^2$-size-bias of $\mu_n$.
As $e^{-sb^2/2}$ strictly decreases
with $b^2$, the reverse covariance
inequality gives $r_{n+1}<r_n$.
Explicitly, the numerator of
$r_n-r_{n+1}$ is one half of
\[
\iint (bc)^{2n}W_s(b)W_s(c)
(b^2-c^2)
\bigl(e^{-sc^2/2}-e^{-sb^2/2}\bigr)
\,db\,dc,
\]
which is strictly positive because
$W_s$ gives positive weight to
different absolute values.

Finally let $0<R<L$. From
\eqref{eq:duhamelbn}--\eqref{eq:duhamelun},
\begin{align}
0<r_n&\le e^{-sR^2/2}
+\frac{\int_{|b|\le R}b^{2n}W_s(b)\,db}
{\int_{\mathbb R}b^{2n}W_s(b)\,db}
\nonumber\\
&\le e^{-sR^2/2}
+\left(\frac RL\right)^{2n}
\frac{\int_{|b|\le R}W_s(b)\,db}
{\int_{|b|\ge L}W_s(b)\,db}.
\label{eq:duhamelrate}
\end{align}
Both integrals in the final quotient
are finite and its denominator is
positive. Letting $n\to\infty$
and then $R\to\infty$ proves
$r_n\to0$. Equation
\eqref{eq:duhamelsaturation}
now follows from $D_n/B_n=1-r_n$.
\end{proof}

\paragraph{Interpretation and an exactly solvable benchmark.}
The theorem is a real obstruction to
a quantitative shortcut, not a refutation
of the original division-free inequality:
at $x=0$ that inequality holds strictly
for each finite order, but its relative
slack must vanish as $n\to\infty$.
For the Gaussian reference measure
$\nu(u)=e^{-au^2}$, $a>0$,
the two variables $w,b$ separate,
and a direct Gaussian moment
calculation gives the exact formula
\begin{equation}
\frac{U_n(s,0)}{B_n(s)}
=\left(\frac{a}{a+s}\right)^{n+1/2}.
\label{eq:duhamelgaussian}
\end{equation}
Thus high-order budget saturation
holds even for a Fourier transform
without any zeros.
The non-circular use of the theta
representation is the explicit
positive correlation density
$W_s$, not a hidden assumption
of all-order Laguerre positivity.

\paragraph{Tests and the remaining gap.}
The rational covariance identities and
Gaussian benchmark are reproduced by
\path{certificates/xi_duhamel_order_saturation_exact.py};
their finite algebraic gates are in
\path{XiKernel/DuhamelOrderSaturation.lean}.
An independent numerical quadrature
of the genuine theta density at
$s=1/4$ gives the values
\[
\begin{array}{c|rrrrr}
n&1&5&20&80&160\\ \hline
D_n/B_n&
0.0083901&0.0270099&0.0739146&
0.1669065&0.2314472
\end{array}
\]
as a diagnostic only, not an interval
certificate. The analytic proof
does not rely on these values.

The division-free source inequality
remains the desired statement for
\emph{all} $(n,s,x)$. Theorem
\ref{thm:duhamelsaturation} shows that
even at $x=0$ one cannot expect any
strict margin uniform in $n$.
At nonzero $x$, the Fourier cosine
factor can cause further cancellations.
A genuine RH proof must therefore
establish a theta-specific \emph{order-sensitive},
sign-preserving Fourier inequality,
or a different global zero-exclusion
principle. Such a theorem has
not yet been established.
